\documentclass[12pt]{article}
\usepackage{notes-project}

\begin{document}

\begin{center}
  {\LARGE\bfseries Lecture 1: Finding Solutions}\\[0.25em]
  {\large Introduction to Data Analysis Techniques}\\[0.15em]
\end{center}

\section{Data Analysis in Science} 

\begin{dfn}[Analytical Method and Analytical Solution]{analytical-method}
  When we use \emph{analytical methods} to solve a problem, we are manipulating the mathematical expressions until unknowns are being isolated into a standard mathematical form. 

  An \emph{analytical solution} is a solution for a problem such that the unknowns are expressed in a closed form. 
\end{dfn}

\begin{eg}[Analytical Solution of a Quadratic Equation]{quadratic-equation}
  Consider the quadratic equation
  \begin{equation}
    ax^2 + bx + c = 0,
  \end{equation}
  where $a$, $b$, and $c$ are constants. The analytic solution for $x$ can be found using the quadratic formula:
  \begin{equation}
    x = \frac{-b \pm \sqrt{b^2 - 4ac}}{2a}.
  \end{equation} 
\end{eg}

\begin{dfn}[Numerical Method and Numerical Solution]{numerical-method}
  \emph{Numerical methods} are methods to obtain numerical solutions. 

  A \emph{numerical solution} is a solution for a problem such that the unknowns are expressed in an numerical form.
\end{dfn} 

\begin{eg}[Numerical Solution of a Trigonometric Equation]{trig-equation}
  Consider the trigonometric equation
  \begin{equation}
    x = \beta \cos x
  \end{equation} 
  where $\beta$ is a constant. 

  It is very difficult to find an analytical solution for this equation using algebraic manipulations. However, we can use numerical methods to find an approximate solution for specific $\beta$ values. 

  The numerical solutions can be 
  \[
  \beta = 1, \quad x \approx 0.73909, \quad \beta = 0.5, \quad x \approx 0.4508 
  \]
\end{eg} 

\begin{rem}[Advantages and Disadvantages of Numerical Methods]{numerical-methods}
  \begin{enumerate}[nosep] 
    \item For numerical methods, we must have all parameters of the problem known. 
    \item If parameters change, solution must be re-done.
    \item Solutions may be not exact. 
    \item Numerical solutions often considered inferior to analytic solutions, if exist. \textit{ie, we should use analytical methods whenever possible or too difficult.}
  \end{enumerate}

\end{rem} 

\section{Using Taylor Series to Approximate Functions} 
\begin{dfn}[Taylor Series]{taylor-series}
  For a function $f(x)$ that is infinitely differentiable at a point $a$, then: 
  \[f(x) = \sum_{n=0}^{\infty} \frac{f^{(n)}(a)}{n!}(x-a)^n\]
\end{dfn}

\begin{rem}[Taylor Series Approximation]{taylor-series-approx}
  It is impossible to compute the infinite sum in practice. When we find the analytical solution of a problem, we can use the first few terms of the \xref[Taylor Series]{def:taylor-series} to approximate the function.

  The more terms we use, the better the approximation.
\end{rem}

\begin{eg}[Electric Field of a Dipole for $x \gg a$]{dipole-field}
  Consider the electric field on the axis of a dipole with charges $+q$ at $x = a$ and $-q$ at $x = -a$:
  \[E = \frac{q}{4\pi\varepsilon_0} \left[
    \frac{1}{(x-a)^2} - \frac{1}{(x+a)^2}
  \right]\]

  Take the Taylor series expansion of $E$ for $x \gg a$:
  \[E = \frac{q}{4\pi\varepsilon_0 \, x^2} \left[
    4\frac{a}{x} + 8\left(\frac{a}{x}\right)^3 + 12\left(\frac{a}{x}\right)^5 + \cdots
  \right]\]

  Keeping the first two terms gives an approximation for $x \gg a$:
  \[E \approx \frac{q}{4\pi\varepsilon_0 \, x^2} \left[
    4\frac{a}{x} + 8\left(\frac{a}{x}\right)^3
  \right]\]

  Which means, for $x \gg a$, $E \propto x^{-3}$.
\end{eg}

\section{Fixed Point Iteration Method} 

\begin{dfn}[Fixed Point Iteration Method]
  The idea for \emph{fixed point iteration method} is to first find a function $g(x)$ such that root $r$ is a \xref[fixed point]{fundamental-computing::lec07::def:fixed-point} of $g(x)$, i.e., $g(r) = r$ 
\end{dfn}

\end{document}
